\documentclass[12pt]{article}

\usepackage[utf8]{inputenc}
\usepackage[spanish]{babel}
\usepackage[shortlabels]{enumitem}

\usepackage{mathtools}
\usepackage{amssymb}
\usepackage{amsthm}

\usepackage{geometry}
\geometry{margin=2.5cm}

\usepackage{hyperref}

\title{Solucion Ejercicios Limites Y Continuidad}
\author{}

\begin{document}
\maketitle

\section*{Solucion Ejercicio 1}

Sea
\[
A=\sqrt{\,1+\ln(1+x)\,},\qquad B=\sqrt{\,1-\ln(1+x^2)\,}.
\]
El limite pedido es
\[
\lim_{x\to 0}\frac{A-B}{x}.
\]

Racionalizamos multiplicando y dividiendo por el conjugado $A+B$:
\[
\frac{A-B}{x}\cdot \frac{A+B}{A+B}
=\frac{A^2-B^2}{x(A+B)}.
\]

Calculamos $A^2-B^2$:
\[
A^2-B^2=\bigl(1+\ln(1+x)\bigr)-\bigl(1-\ln(1+x^2)\bigr)
=\ln(1+x)+\ln(1+x^2).
\]

Por tanto,
\[
\frac{\sqrt{\,1+\ln(1+x)\,}-\sqrt{\,1-\ln(1+x^2)\,}}{x}
=\frac{\ln(1+x)+\ln(1+x^2)}{x\left(\sqrt{\,1+\ln(1+x)\,}+\sqrt{\,1-\ln(1+x^2)\,}\right)}.
\]

Separamos:
\[
\frac{\ln(1+x)+\ln(1+x^2)}{x\left(\sqrt{\,1+\ln(1+x)\,}+\sqrt{\,1-\ln(1+x^2)\,}\right)} = 
\]
\[
\frac{1}{\left(\sqrt{\,1+\ln(1+x)\,}+\sqrt{\,1-\ln(1+x^2)\,}\right)} \cdot \frac{\ln(1+x)+\ln(1+x^2)}{x}
\]
Por lo que:

\[
\lim_{x\to 0}
\frac{\ln(1+x)+\ln(1+x^2)}{x\left(\sqrt{\,1+\ln(1+x)\,}+\sqrt{\,1-\ln(1+x^2)\,}\right)}
=
\frac{1}{2}
\lim_{x\to 0}\frac{\ln(1+x)+\ln(1+x^2)}{x}.
\]

Separamos:
\[
\lim_{x\to 0}\frac{\ln(1+x)+\ln(1+x^2)}{x}
=
\lim_{x\to 0}\frac{\ln(1+x)}{x}
+
\lim_{x\to 0}\frac{\ln(1+x^2)}{x}.
\]

Primer termino (aplicamos infinitesimos):
\[
\lim_{x\to 0}\frac{\ln(1+x)}{x}\simeq\lim_{x\to 0}\frac{x}{x} = 1.
\]

Segundo termino (aplicamos infinitesimos):
\[
\lim_{x\to 0}\frac{\ln(1+x^2)}{x} \simeq \lim_{x\to 0} \frac{x^2}{x} = 0
\]

Por tanto,
\[
\lim_{x\to 0}\frac{\ln(1+x)+\ln(1+x^2)}{x}=1.
\]

Finalmente,
\[
\lim_{x\to 0}\frac{\sqrt{\,1+\ln(1+x)\,}-\sqrt{\,1-\ln(1+x^2)\,}}{x}
=\frac{1}{2}.
\]

\section*{Solucion Ejercicio 2}

Queremos calcular
\[
\lim_{x\to 0}\frac{1+2(\cos x)^3-3\sqrt{\cos(2x)}}{(\sin x)^2}.
\]
Sea
\[
N(x)=1+2(\cos x)^3-3\sqrt{\cos(2x)}.
\]
Entonces el limite es
\[
\lim_{x\to 0}\frac{N(x)}{\sin^2 x}.
\]

Aplicamos infinitesimos $\sin x \sim x$
\[
\frac{N(x)}{\sin^2 x}
\simeq \frac{N(x)}{x^2}
\]


\subsection*{2) Calcular $\displaystyle \lim_{x\to 0}\frac{N(x)}{x^2}$ con L'Hopital (dos veces)}
Primero notamos que $N(0)=1+2\cdot 1-3\cdot 1=0$, asi que
\[
\lim_{x\to 0}\frac{N(x)}{x^2}
\ \text{es una forma}\ \frac{0}{0}.
\]
Aplicamos L'Hopital:
\[
\lim_{x\to 0}\frac{N(x)}{x^2}
=\lim_{x\to 0}\frac{N'(x)}{2x}.
\]
Vemos que tambien es $\frac{0}{0}$ (lo veremos al evaluar $N'(0)$), asi que aplicamos L'Hopital otra vez:
\[
\lim_{x\to 0}\frac{N'(x)}{2x}
=\lim_{x\to 0}\frac{N''(x)}{2},
\]

Ahora calculamos derivadas.

\subsection*{3) Derivadas de $N(x)$}
\[
N(x)=1+2(\cos x)^3-3\sqrt{\cos(2x)}.
\]

\paragraph{Primera derivada.}
\[
\frac{d}{dx}\bigl(2(\cos x)^3\bigr)=2\cdot 3(\cos x)^2(-\sin x)=-6\cos^2 x\,\sin x.
\]
Para la raiz, usando regla de la cadena:
\[
\frac{d}{dx}\sqrt{\cos(2x)}
=\frac{1}{2\sqrt{\cos(2x)}}\cdot(-\sin(2x)\cdot 2)
=-\frac{\sin(2x)}{\sqrt{\cos(2x)}}.
\]
Luego
\[
\frac{d}{dx}\bigl(-3\sqrt{\cos(2x)}\bigr)
=-3\left(-\frac{\sin(2x)}{\sqrt{\cos(2x)}}\right)
=3\frac{\sin(2x)}{\sqrt{\cos(2x)}}.
\]
Por tanto,
\[
N'(x)=-6\cos^2 x\,\sin x+3\frac{\sin(2x)}{\sqrt{\cos(2x)}}.
\]
Evaluando en $x=0$:
\[
N'(0)=-6\cdot 1\cdot 0+3\cdot \frac{0}{1}=0,
\]
por eso era otra vez forma $\frac{0}{0}$.

\paragraph{Segunda derivada.}
Derivamos cada termino de $N'(x)$.

\textbf{(i) Derivada de $-6\cos^2 x\,\sin x$.}
\[
\frac{d}{dx}\bigl(\cos^2 x\,\sin x\bigr)
=(\cos^2 x)'\sin x+\cos^2 x(\sin x)'
= (2\cos x(-\sin x))\sin x+\cos^2 x\cos x.
\]
Asi
\[
(\cos^2 x\,\sin x)'=-2\cos x\sin^2 x+\cos^3 x,
\]
y por tanto
\[
\frac{d}{dx}\bigl(-6\cos^2 x\,\sin x\bigr)
=-6\bigl(-2\cos x\sin^2 x+\cos^3 x\bigr)
=12\cos x\sin^2 x-6\cos^3 x.
\]

\textbf{(ii) Derivada de $3\dfrac{\sin(2x)}{\sqrt{\cos(2x)}}$.}
Sea
\[
f(x)=\sin(2x),\qquad g(x)=\sqrt{\cos(2x)}.
\]
Entonces
\[
\left(\frac{f}{g}\right)'=\frac{f'g-fg'}{g^2}.
\]
Calculamos:
\[
f'(x)=2\cos(2x),
\qquad
g'(x)=-\frac{\sin(2x)}{\sqrt{\cos(2x)}}=-\frac{f(x)}{g(x)}.
\]
Sustituyendo:
\[
\left(\frac{f}{g}\right)'
=\frac{2\cos(2x)\,g-f\left(-\frac{f}{g}\right)}{g^2}
=\frac{2\cos(2x)\,g+\frac{f^2}{g}}{g^2}
=\frac{2\cos(2x)}{g}+\frac{f^2}{g^3}.
\]
Volviendo a $g=\sqrt{\cos(2x)}$ y $f=\sin(2x)$:
\[
\left(\frac{\sin(2x)}{\sqrt{\cos(2x)}}\right)'
=\frac{2\cos(2x)}{\sqrt{\cos(2x)}}
+\frac{\sin^2(2x)}{(\cos(2x))^{3/2}}.
\]
Multiplicando por $3$:
\[
\frac{d}{dx}\left(3\frac{\sin(2x)}{\sqrt{\cos(2x)}}\right)
=3\left(
\frac{2\cos(2x)}{\sqrt{\cos(2x)}}
+\frac{\sin^2(2x)}{(\cos(2x))^{3/2}}
\right).
\]

\paragraph{Conclusion para $N''(x)$.}
\[
N''(x)=\bigl(12\cos x\sin^2 x-6\cos^3 x\bigr)
+3\left(
\frac{2\cos(2x)}{\sqrt{\cos(2x)}}
+\frac{\sin^2(2x)}{(\cos(2x))^{3/2}}
\right).
\]

Evaluamos en $x=0$:
\[
\cos 0=1,\ \sin 0=0
\]
entonces
\[
N''(0)=\left(12\cdot 1\cdot 0-6\cdot 1\right)
+3\left(\frac{2\cdot 1}{1}+\frac{0}{1}\right)
=-6+6=0.
\]

\subsection*{4) Volver al limite}
\[
\lim_{x\to 0}\frac{N''(x)}{2}=\frac{N''(0)}{2}=0.
\]
Por L'Hopital dos veces:
\[
\lim_{x\to 0}\frac{N(x)}{x^2}=0.
\]

\[
\boxed{
\lim_{x\to 0}\frac{1+2(\cos x)^3-3\sqrt{\cos(2x)}}{(\sin x)^2}=0.
}
\]

\section*{Solucion Ejercicio 3}

La funcion esta definida en el dominio explicito $(-\infty,2]$ mediante dos expresiones:
\[
f(x)=
\begin{cases}
\dfrac{a}{1+e^{-bx}}, & x<0,\\[6pt]
\dfrac{3x+2}{x^2-x-6}, & 0\le x\le 2.
\end{cases}
\]
Para que $f$ sea continua en todo su dominio, debe ser continua:
\begin{itemize}
\item en $(-\infty,0)$,
\item en $[0,2]$,
\item y ademas en el punto crítico $x=0$.
\end{itemize}

\subsection*{1) Continuidad en $(-\infty,0)$}
Para $x<0$:
\[
f(x)=\frac{a}{1+e^{-bx}}.
\]
La funcion exponencial es continua y $1+e^{-bx}>0$ para todo $x$, asi que no hay division por cero.
Por tanto, esta rama es continua en todo $(-\infty,0)$ para cualesquiera $a,b\in\mathbb{R}$.

\subsection*{2) Continuidad en $[0,2]$}
Para $0\le x\le 2$:
\[
f(x)=\frac{3x+2}{x^2-x-6}.
\]
Una funcion racional es continua donde el denominador no se anula. Factorizamos:
\[
x^2-x-6=(x-3)(x+2).
\]
Las posibles discontinuidades estarian en $x=3$ y $x=-2$, pero ninguna pertenece al intervalo $[0,2]$.
Luego, en $[0,2]$ el denominador nunca es cero y esta rama es continua en todo $[0,2]$.

\subsection*{3) Continuidad en el punto crítico $x=0$}
Como $0$ pertenece al segundo tramo ($[0,2]$), el valor de la funcion en $0$ es:
\[
f(0)=\frac{3\cdot 0+2}{0^2-0-6}=\frac{2}{-6}=-\frac{1}{3}.
\]

Para que $f$ sea continua en $x=0$, debe cumplirse:
\[
\lim_{x\to 0^-} f(x)=f(0).
\]
Calculamos el limite por la izquierda usando el primer tramo:
\[
\lim_{x\to 0^-}\frac{a}{1+e^{-bx}}.
\]
\[
\lim_{x\to 0^-}\frac{a}{1+e^{-bx}}
=\frac{a}{1+1}=\frac{a}{2}.
\]
Imponiendo continuidad:
\[
\frac{a}{2}=-\frac{1}{3}
\quad\Longrightarrow\quad
a=-\frac{2}{3}.
\]

\subsection*{4) Conclusion}
La condicion de continuidad en $x=0$ obliga a
\[
a=-\frac{2}{3},
\]
y el parametro $b$ no queda restringido.

\[
\boxed{\text{La funcion es continua en todo }(-\infty,2]\text{ si y solo si } a=-\frac{2}{3}\text{ y } b\in\mathbb{R}.}
\]

\section*{Solucion Ejercicio 4}

Antes de empezar, notamos que para $x\ne 0$ aparece el termino $1/\sen(x)$, luego la expresion
\[
\frac{\sen(x)}{1+e^{1/\sen(x)}}
\]
solo tiene sentido cuando $\sen(x)\ne 0$. En particular, en $x=\pi$ (donde $\sen(\pi)=0$) la formula no esta definida.

\subsection*{a) Continuidad en $x=0$}

Por definicion,
\[
f(0)=0.
\]
Para estudiar continuidad, calculamos el limite cuando $x\to 0$.

\[
\lim_{x\to 0} \frac{\sen(x)}{1+e^{1/\sen(x)}} = 0.
\]
Y como $f(0)=0$, resulta
\[
\lim_{x\to 0} f(x)=f(0),
\]
luego $f$ es continua en $x=0$.

\subsection*{b) Continuidad en $x=\pi/2$}

Como $\sen(\pi/2)=1\ne 0$, la expresion esta definida en un entorno de $\pi/2$.
Ademas, las funciones $\sen(x)$ y $e^{1/\sen(x)}$ son continuas donde $\sen(x)\ne 0$,
y el denominador $1+e^{1/\sen(x)}$ es siempre positivo, asi que no se anula.

Por tanto, $f$ es continua en $x=\pi/2$ y su valor es
\[
f\left(\frac{\pi}{2}\right)=\frac{1}{1+e^{1}}=\frac{1}{1+e}.
\]

\subsection*{c) Continuidad en $x=\pi$}

En $x=\pi$ se tiene $\sen(\pi)=0$, por lo que la expresion
\[
\frac{\sen(x)}{1+e^{1/\sen(x)}}
\]
no esta definida en $x=\pi$ (porque aparece $1/\sen(\pi)$). Por tanto, tal y como esta dada la definicion, $f(\pi)$ no existe y $f$ no puede ser continua en $x=\pi$.

Aun asi, podemos estudiar el limite cuando $x\to \pi$ para clasificar la discontinuidad.

\[
\lim_{x\to \pi}\frac{\sen(x)}{1+e^{1/\sen(x)}} = 0
\]

\paragraph{Conclusion en $x=\pi$.}
$\lim_{x\to \pi}f(x)$ existe, pero $f(\pi)$ no esta definida con la definicion dada, luego hay una \textbf{discontinuidad evitable} en $x=\pi$.

De hecho, si se define $f(\pi)=0$, la funcion quedaria continua en $x=\pi$.

\subsection*{Resumen}
\[
\boxed{
\text{$f$ es continua en $x=0$ y en $x=\pi/2$. En $x=\pi$ tiene discontinuidad evitable.}
}
\]

\section*{Solucion Ejercicio 5}

Consideramos
\[
L=\lim_{x\to 0}\frac{\sqrt{1+x\sen(x)}-\sqrt{\cos(2x)}}{\tan^2\!\left(\frac{x}{2}\right)}.
\]

\subsection*{1) Forma indeterminada}

Cuando $x\to 0$:
\[
x\sen(x)\to 0, \qquad \cos(2x)\to 1, \qquad \tan\!\left(\frac{x}{2}\right)\to 0.
\]
Por tanto,
\[
\sqrt{1+x\sen(x)}\to 1, \qquad \sqrt{\cos(2x)}\to 1,
\]
y el limite es de tipo
\[
\frac{0}{0}.
\]

\subsection*{2) Racionalizar el numerador}

Multiplicamos y dividimos por el conjugado:
\[
\frac{\sqrt{1+x\sen(x)}-\sqrt{\cos(2x)}}{\tan^2\!\left(\frac{x}{2}\right)}
\cdot
\frac{\sqrt{1+x\sen(x)}+\sqrt{\cos(2x)}}{\sqrt{1+x\sen(x)}+\sqrt{\cos(2x)}}.
\]
El numerador queda
\[
(1+x\sen(x))-\cos(2x).
\]
Luego
\[
L=\lim_{x\to 0}
\frac{1+x\sen(x)-\cos(2x)}
{\tan^2\!\left(\frac{x}{2}\right)\left(\sqrt{1+x\sen(x)}+\sqrt{\cos(2x)}\right)}.
\]

\subsection*{3) Simplificar usando relaciones trigonometricas}

Usamos
\[
1-\cos(2x)=2\sen^2 x.
\]
Entonces
\[
1+x\sen(x)-\cos(2x)
=2\sen^2 x + x\sen(x)
=\sen(x)\bigl(2\sen(x)+x\bigr).
\]
Por tanto,
\[
L=\lim_{x\to 0}
\frac{\sen(x)\bigl(2\sen(x)+x\bigr)}
{\tan^2\!\left(\frac{x}{2}\right)\left(\sqrt{1+x\sen(x)}+\sqrt{\cos(2x)}\right)}.
\]

Recordamos
\[
\tan\!\left(\frac{x}{2}\right)
=\frac{\sen\!\left(\frac{x}{2}\right)}{\cos\!\left(\frac{x}{2}\right)}
\quad\Rightarrow\quad
\tan^2\!\left(\frac{x}{2}\right)
=\frac{\sen^2\!\left(\frac{x}{2}\right)}{\cos^2\!\left(\frac{x}{2}\right)}.
\]
Sustituyendo,
\[
L=\lim_{x\to 0}
\frac{\sen(x)\bigl(2\sen(x)+x\bigr)\cos^2\!\left(\frac{x}{2}\right)}
{\sen^2\!\left(\frac{x}{2}\right)\left(\sqrt{1+x\sen(x)}+\sqrt{\cos(2x)}\right)}.
\]

Usamos ahora
\[
\sen(x)=2\sen\!\left(\frac{x}{2}\right)\cos\!\left(\frac{x}{2}\right).
\]
Entonces
\[
\sen(x)\cos^2\!\left(\frac{x}{2}\right)
=
2\sen\!\left(\frac{x}{2}\right)\cos^3\!\left(\frac{x}{2}\right).
\]
Sustituyendo:
\[
L=\lim_{x\to 0}
\frac{2\sen\!\left(\frac{x}{2}\right)\cos^3\!\left(\frac{x}{2}\right)\bigl(2\sen(x)+x\bigr)}
{\sen^2\!\left(\frac{x}{2}\right)\left(\sqrt{1+x\sen(x)}+\sqrt{\cos(2x)}\right)}.
\]
Simplificamos un $\sen(x/2)$:
\[
L=\lim_{x\to 0}
\frac{2\cos^3\!\left(\frac{x}{2}\right)\bigl(2\sen(x)+x\bigr)}
{\sen\!\left(\frac{x}{2}\right)\left(\sqrt{1+x\sen(x)}+\sqrt{\cos(2x)}\right)}.
\]

\subsection*{4) Evaluar limites}

Cuando $x\to 0$:
\[
\cos\!\left(\frac{x}{2}\right)\to 1,
\qquad
\sqrt{1+x\sen(x)}\to 1,
\qquad
\sqrt{\cos(2x)}\to 1.
\]
Luego
\[
\sqrt{1+x\sen(x)}+\sqrt{\cos(2x)}\to 2.
\]

Para el factor restante,
\[
\frac{2\sen(x)+x}{\sen\!\left(\frac{x}{2}\right)}
=
\frac{2\sen(x)}{\sen\!\left(\frac{x}{2}\right)}
+\frac{x}{\sen\!\left(\frac{x}{2}\right)}.
\]
Usando $\sen(x)=2\sen(x/2)\cos(x/2)$:
\[
\frac{2\sen(x)}{\sen(x/2)}
=4\cos\!\left(\frac{x}{2}\right)\to 4.
\]
Ademas,
\[
\frac{x}{\sen(x/2)} \simeq \frac{x}{x/2}\to 2.
\]
Por tanto,
\[
\frac{2\sen(x)+x}{\sen(x/2)}\to 6.
\]

Finalmente,
\[
L=\frac{2\cdot 1^3\cdot 6}{2}=6.
\]

\[
\boxed{6}
\]
\section*{Solucion Ejercicio 6}

Sea
\[
L=\lim_{x\to \pi/2}\left(\frac{e^x\cos(x)+1}{e^{\cos(x)}}\right)^{1/\cos(x)}.
\]
Cuando $x\to \pi/2$ se tiene $\cos(x)\to 0$, luego la base tiende a $1$ y el exponente tiende a infinito. Es una indeterminacion del tipo $1^\infty$.

\subsection*{1) Pasar a logaritmos}
Tomamos logaritmo:
\[
\ln L=\lim_{x\to \pi/2}\frac{1}{\cos(x)}\ln\left(\frac{e^x\cos(x)+1}{e^{\cos(x)}}\right),
\]
siempre que el limite de la derecha exista. Dentro del logaritmo:
\[
\ln\left(\frac{e^x\cos(x)+1}{e^{\cos(x)}}\right)
=\ln\bigl(1+e^x\cos(x)\bigr)-\ln\bigl(e^{\cos(x)}\bigr)
=\ln\bigl(1+e^x\cos(x)\bigr)-\cos(x).
\]
Por tanto,
\[
\ln L=\lim_{x\to \pi/2}\left(\frac{\ln\bigl(1+e^x\cos(x)\bigr)}{\cos(x)}-1\right).
\]

\subsection*{2) Usar el infinitesimo $\displaystyle \lim_{x\to 0}\ln(1+f(x)) \sim f(x)$}
Escribimos
\[
\frac{\ln\bigl(1+e^x\cos(x)\bigr)}{\cos(x)} \simeq \frac{e^x\cos(x)}{\cos(x)} = e^x
\]
Por tanto,
\[
\lim_{x\to \pi/2}\frac{\ln\bigl(1+e^x\cos(x)\bigr)}{\cos(x)} = e^{\pi/2}.
\]
Sustituyendo en la expresion de $\ln L$:
\[
\ln L=\left(e^{\pi/2}\right)-1=e^{\pi/2}-1.
\]

\subsection*{3) Volver de logaritmos}
Como $\ln L=e^{\pi/2}-1$, se obtiene
\[
L=e^{\,e^{\pi/2}-1}.
\]

\[
\boxed{\displaystyle \lim_{x\to \pi/2}\left(\frac{e^x\cos(x)+1}{e^{\cos(x)}}\right)^{1/\cos(x)}
= e^{\,e^{\pi/2}-1}}
\]

\section*{Solucion Ejercicio 7}

Para que $f$ sea continua en $x=0$ se debe cumplir:
\[
\alpha=f(0)=\lim_{x\to 0}\frac{\ln\left(\frac{\sen(x)}{x}\right)}{x^2},
\]
siempre que ese limite exista.

Sea
\[
L=\lim_{x\to 0}\frac{\ln\left(\frac{\sen(x)}{x}\right)}{x^2}.
\]
El limite tiene forma $\frac{0}{0}$.

Vamos a plantear el numerador de la forma $\ln(f(x)) \sim \ln(1 + u) \sim u$ tal que $1+u = f(x) \rightarrow u = f(x) - 1$

Nota: Este cambio no vale para todos los casos. Tienen que ser funciones equivalentes en el punto de trabajo

Por lo que:
\[
L=\lim_{x\to 0}\frac{\frac{\sen(x)}{x} - 1}{x^2} = \lim_{x\to 0}\frac{\sen(x)-x}{x^3}
\]

Aplicamos L'Hopital dos veces
\[
\lim_{x\to 0}\frac{\sen(x)-x}{x^3}=-\frac{1}{6}.
\]
Por tanto,
\[
\lim_{x\to 0}\frac{\ln\left(\frac{\sen(x)}{x}\right)}{x^2}-\frac{1}{6}.
\]

Para continuidad en $x=0$ necesitamos $\alpha=L$, asi que
\[
\boxed{\alpha=-\frac{1}{6}}.
\]

\section*{Solucion Ejercicio 8}

Para que $f$ sea continua en $x=0$ debe cumplirse:
\[
\lim_{x\to 0^-} f(x)=f(0)=\lim_{x\to 0^+} f(x).
\]

\subsection*{1) Valor de la funcion en $x=0$}

Como $0$ pertenece al segundo tramo,
\[
f(0)=\frac{2\cdot 0+1}{0-1}=-1.
\]

\subsection*{2) Limite por la derecha ($x\to 0^+$)}

Para $x>0$ se usa la segunda expresion, que es una funcion racional continua donde el denominador no se anula. Luego
\[
\lim_{x\to 0^+} f(x)=f(0)=-1.
\]

\subsection*{3) Limite por la izquierda ($x\to 0^-$)}

Para $x<0$:
\[
f(x)=\frac{\sqrt{1+kx}-\sqrt{1-kx}}{x}.
\]
Sustituyendo directamente $x=0$ se obtiene forma $\frac{0}{0}$, asi que racionalizamos multiplicando por el conjugado:
\[
\frac{\sqrt{1+kx}-\sqrt{1-kx}}{x}
\cdot
\frac{\sqrt{1+kx}+\sqrt{1-kx}}{\sqrt{1+kx}+\sqrt{1-kx}}.
\]
El numerador queda
\[
(1+kx)-(1-kx)=2kx.
\]
Por tanto,
\[
f(x)=\frac{2kx}{x\bigl(\sqrt{1+kx}+\sqrt{1-kx}\bigr)}
=\frac{2k}{\sqrt{1+kx}+\sqrt{1-kx}}.
\]
Ahora el limite es inmediato:
\[
\lim_{x\to 0^-} f(x)
=\frac{2k}{1+1}=k.
\]

\subsection*{4) Imponer continuidad}

Para continuidad en $x=0$:
\[
\lim_{x\to 0^-} f(x)=f(0).
\]
Es decir,
\[
k=-1.
\]

\subsection*{5) Conclusion}

\[
\boxed{\text{La funcion es continua en } x=0 \text{ si y solo si } k=-1.}
\]
\end{document}